\section{适用范围、归属与限制}\label{sec:scope}

本文结论对普通的任意六维 Kakeya 集合成立的含义，是在明列四维有限输入及
Katz--Tao 具名递推下的推论。作为完整核验责任，读者需同时检查四维有限接口的原证明与所引高级输入；
作者署名、档案哈希、数值试验或内部审查标签都不是正确性证明。

需要特别区分以下三个层次：四维集合的满维数定理；
适用于带标签原标记的四维有限方向质量估计；
由该估计与递推得到的高维后果。本文使用第二层作为真实输入。
其输出对每次新的有限族统一成立，不要求这些族来自同一跨尺度集合。
但它有逐线长标记门槛，不能据此声称任意小着色密度的尖锐 Kakeya 极大函数估计。

方向删选保留原角上界，是未归一化的子族上界；若改为归一化概率律，必须支付剩余质量分母。
平均着色下界先经过重线删选，再进入逐线输入，故公式\eqref{eq:finite-sharp-budget}确有两个 $\lambda$。
投影后只有在碰撞图删选之后才有分离方向，不能免费继承投影 Frostman 控制。
任意环境空间和实角参数是递推的实际需要，并非对结论作形式换名。

相对\cite{KT,HRZ,GLX}中明确列出的普通六维 Hausdorff 基线，
式\eqref{eq:comparison}给出固定正指数改善。
四维原稿\cite[Section 11.2]{OAI4}明确记录的是向更高环境投影所给的下界 $4$。
截至前述核验截止日的有界比较没有识别同一 $17/4$ 的明确报告或可具体采用的同类型支配结果，
但有限检索不证明绝对首创、最新最优或不存在隐含推论。

本文不提出新的 Katz--Tao 递推，也不独立占有四维原稿或其高级输入的贡献。
五维满 Hausdorff 维数、一般维度满维数及真实反例问题均未由本文解决。
本稿适合进一步的外部数学审阅；内部交叉审查与外部同行评审严格区分。
